\chapter{Introduction}
\label{ch:introduction}

\section{The Verification Gap in Process Mining Implementations}
\label{sec:intro-gap}

Process mining has matured into a discipline with well-understood algorithms,
published standards, and production-grade open-source libraries. The Alpha
Miner~\cite{vanderaalst2004workflow}, the Inductive
Miner~\cite{leemans2013discovering}, alignment-based conformance
checking~\cite{adriansyah2011conformance}, and the OCEL~2.0
standard~\cite{ocel20_2023} collectively represent decades of theoretical
development. Libraries such as pm4py~\cite{berti2019pm4py} and toolkits built
on top of them bring these algorithms within reach of practitioners writing
ordinary Python or Rust code.

Yet a persistent gap separates the guarantees stated in the literature from
the guarantees that can be verified in a running implementation. A Petri net
soundness theorem~\cite{vanderaalst1998application} holds unconditionally for
the mathematical object it describes. A software representation of that net
holds its invariants only as long as every caller, at every call site, respects
the conventions documented in comments, enforced by runtime assertions, or
validated by tests. None of these mechanisms is structural: none of them
prevents a caller from constructing an unsound net, admitting it as sound, and
passing it downstream without the compiler raising an objection.

This is the \emph{verification gap}: the space between what a type signature
claims and what the compiler is willing to enforce. Closing the gap requires
making the invariant \emph{visible to the type system} --- encoding it as a
type-level distinction that the compiler checks at every use site, without
exception and without the need for runtime evidence.

The \texttt{wasm4pm-compat} crate is an attempt to close this gap for the core
abstractions of object-centric process mining. Its central invariant is a
typed, one-way lifecycle: raw data must be admitted before it can be projected,
exported, or receipted, and admission is only possible through a sanctioned
path that names a specific structural law. The compiler enforces this invariant
unconditionally. No escape hatch exists short of \texttt{unsafe} code, and
the crate unconditionally declares \texttt{\#![forbid(unsafe\_code)]}.

\section{The Central Claim: Type Law as Compile-Time Proof Obligation}
\label{sec:intro-claim}

The central claim of this work is stated precisely:

\begin{quote}
  A \emph{type law} is a structural invariant of a domain that can be
  encoded as a well-typed program such that any violation causes a
  compile-time error. A \emph{compile-time proof obligation} is a type-law
  encoding together with a trybuild fixture that demonstrates the obligation is
  both enforced (compile-fail path exists and fails for the named law) and
  satisfiable (compile-pass path exists and compiles successfully).
\end{quote}

The key safety property is:

\begin{equation}
  \label{eq:type-law-safety}
  \forall T, S_1, S_2, W.\
  \Evidence{T}{S_1}{W} \not\leq \Evidence{T}{S_2}{W}
  \text{ when } S_1 \neq S_2
\end{equation}

That is, no subtyping, coercion, or implicit conversion relates two
\Evidence{}{}{} carriers that differ only in their state parameter. The Rust
type system enforces this through nominal typing: \texttt{Evidence<T, Raw, W>}
and \texttt{Evidence<T, Admitted, W>} are distinct types with no common
supertype, no \texttt{From}/\texttt{Into} implementation bridging them, and no
public constructor on the \texttt{Admitted} variant reachable from outside the
crate. The \emph{only} public path to admitted evidence is through
\texttt{Admit::admit()}, which returns \texttt{Result<Admission<\ldots>,
Refusal<R, W>>} where the reason type \texttt{R} must name a specific
structural law.

This claim is not merely theoretical. The crate's ALIVE gate
(\S\ref{sec:intro-contributions}) provides machine-checked evidence that the
claim holds for every law surface in the codebase.

\section{The Manufacturing Doctrine}
\label{sec:intro-doctrine}

The \texttt{wasm4pm-compat} development model is governed by the
\emph{manufacturing doctrine}: a discipline that treats code commits as
receipts, compile-fail fixtures as proof obligations, and the ALIVE gate as the
certification surface.

Under the manufacturing doctrine:

\begin{itemize}
  \item Every commit that introduces or strengthens a type-law surface must
        also introduce or update a trybuild fixture that proves the law is
        enforced. A commit that adds a type restriction without a corresponding
        compile-fail fixture is not a valid receipt.
  \item A compile-fail fixture that fails for the wrong reason --- a missing
        import, a typo, an unstable feature regression --- is not a valid
        type-law receipt. The fixture must fail for the \emph{intended named
        law}, as verified by the expected \texttt{.stderr} diagnostic.
  \item Every lawful path must have a corresponding compile-pass fixture. A
        type restriction that also blocks the lawful path is a defect, not a
        feature.
  \item The graduation boundary to the \texttt{wasm4pm} execution engine is
        enforced by the \texttt{wasm4pm} Cargo feature. Types that bridge into
        the execution engine are only reachable when the feature is enabled.
        This ensures that the compatibility crate does not accidentally import
        execution-engine concerns into the structure-only layer.
\end{itemize}

The doctrine is not a testing strategy. It is a manufacturing strategy: the
commit history is itself an artifact, and each commit either advances the
certification state or it does not. Reaching
\textsc{PaperLaw Crown Alive 004} required 601~receipt-bearing commits
manufactured across May~2026, each traceable to a specific paper, standard, or
law in the process mining literature.

\section{Contributions}
\label{sec:intro-contributions}

This work makes the following contributions:

\begin{itemize}
  \item \textbf{Evidence lifecycle type.} The \Evidence{T}{State}{W} type
        encodes a one-way process-evidence lifecycle at the type level, with
        zero runtime cost. State tokens and witness markers are
        \texttt{PhantomData} fields erased at compile time.

  \item \textbf{Witness-indexed admission.} The \texttt{Admit} trait couples
        admission to a specific named witness (OCEL~2.0, XES~1849, WF-net
        soundness, etc.), making it impossible to confuse evidence admitted
        under one standard with evidence admitted under another.

  \item \textbf{Named-law refusal.} Every refusal carries a specific named law
        as its reason type. String-typed or catch-all refusals are rejected at
        compile time by the type system. This ensures that diagnostic
        information is structured and machine-readable.

  \item \textbf{Loss accounting protocol.} The \texttt{Project} trait enforces
        a \texttt{LossPolicy} decision before any lossy transformation occurs,
        and requires a \texttt{LossReport} on every non-refusing path. Silent
        structure loss is a compile-time defect.

  \item \textbf{Const-generic law surfaces.} \texttt{WfNetConst<SOUNDNESS>},
        \texttt{ConditionCell<BITS>}, \texttt{Between01<NUM, DEN>}, and
        \texttt{TypedLoopNode<ARITY>} encode quantitative invariants directly
        in const generic parameters, checked by the compiler via
        \texttt{Require}/\texttt{IsTrue} bounds.

  \item \textbf{The ALIVE gate.} A trybuild-based certification surface
        consisting of compile-fail and compile-pass fixtures. Reaching
        \textsc{PaperLaw Crown Alive 004}: 196~compile-fail receipts,
        406~compile-pass receipts, 37~modules, 98~papers.

  \item \textbf{The manufacturing doctrine.} A development discipline in
        which commits are receipts, fixtures are proof obligations, and the
        graduation boundary to the execution engine is a type contract. The
        doctrine is the subject of Chapter~\ref{ch:doctrine}.
\end{itemize}

\section{A Motivating Example}
\label{sec:intro-example}

To make the type-law approach concrete, consider the simplest possible misuse
of a process-evidence type: treating raw, unvalidated data as if it had already
been admitted against a standard.

In a runtime-checked library, this error is latent. The raw data may carry a
boolean field \texttt{is\_admitted = false}, but nothing prevents a caller from
passing the object to a function that expects admitted evidence and only
checking the field at runtime --- or not checking it at all.

In \texttt{wasm4pm-compat}, the same error is a compile-time failure:

\begin{lstlisting}[language=Rust, caption={Type-level lifecycle enforcement.
  \texttt{Evidence<T, Raw, W>} and \texttt{Evidence<T, Admitted, W>} are
  distinct types. The compiler rejects any attempt to use one as the other.},
  label={lst:intro-lifecycle}]
use wasm4pm_compat::evidence::Evidence;
use wasm4pm_compat::state::{Raw, Admitted};
use wasm4pm_compat::witness::Ocel20;

// Raw evidence: the standard has not been verified.
let raw: Evidence<OcelLog, Raw, Ocel20> = Evidence::raw(log);

// The only sanctioned path to admitted evidence:
let admitted: Evidence<OcelLog, Admitted, Ocel20> =
    Ocel20Admitter::admit(raw)?;

// This does not compile. Raw and Admitted are distinct types.
// let wrong: Evidence<OcelLog, Admitted, Ocel20> = raw;
//            ^^^^^^^^^^^^^^^^^^^^^^^^^^^^^^^^^^^^^^^^^^^^^^
// error[E0308]: mismatched types
//   expected `Evidence<OcelLog, Admitted, Ocel20>`
//      found `Evidence<OcelLog, Raw, Ocel20>`
\end{lstlisting}

The type-level distinction between \texttt{Raw} and \texttt{Admitted} is the
same kind of distinction that separates an untrusted network buffer from a
validated domain object in a security-critical system --- except that in
\texttt{wasm4pm-compat}, the enforcer is the compiler rather than a runtime
guard. The state token \texttt{Admitted} is a zero-sized type; it contributes
no bytes to the memory layout of \texttt{Evidence}. The entire enforcement cost
is paid at compile time and then erased.

The same pattern extends across the full lifecycle:

\begin{lstlisting}[language=Rust, caption={Full lifecycle from \texttt{Raw} to
  \texttt{Receipted}. Each arrow is a type transition. None can be reversed or
  skipped without a compiler error.}, label={lst:intro-full-lifecycle}]
// Raw -> Admitted (via Admit::admit)
let admitted = Ocel20Admitter::admit(raw)?;

// Admitted -> Projected (via Project::project)
let projected = OcelProjector::project(
    admitted,
    ProjectionName("flatten_to_xes"),
    AllowNamedProjection,
)?;

// Admitted -> Receipted (via receipt issuance)
let receipted = receipt::issue(admitted, receipt_key);

// Admitted -> Refused (terminal; carries named law)
// Only reachable through admit(), never from outside.
\end{lstlisting}

Crucially, there is no \texttt{Raw -> Receipted} path. There is no
\texttt{Refused -> Admitted} path. The type system enforces a
directed acyclic lifecycle: once evidence is refused, it cannot be
rehabilitated. The only way to obtain admitted evidence is to start from raw
evidence and pass through \texttt{Admit::admit()} with a conforming payload.

\section{Paper Structure}
\label{sec:intro-structure}

The remainder of this book is organized as follows.

Chapter~\ref{ch:background} provides background on process mining abstractions
(event logs, Petri nets, workflow nets, OCEL~2.0, conformance checking) and on
the nightly Rust type system features that make the type-law encoding possible
(\texttt{generic\_const\_exprs}, \texttt{adt\_const\_params},
\texttt{const\_trait\_impl}, \texttt{min\_specialization}).

Chapter~\ref{ch:architecture} describes the three-layer type system in detail:
state tokens, witness markers, and the \Evidence{}{}{} carrier. It documents
every canon module and the invariants each module enforces.

Chapter~\ref{ch:law} presents the type-law kernel: \texttt{law.rs},
\texttt{petri.rs}, \texttt{conformance.rs}, \texttt{process\_tree.rs},
\texttt{powl.rs}, and the const-generic bound machinery
(\texttt{Assert}/\texttt{IsTrue}/\texttt{Require}).

Chapter~\ref{ch:alive} describes the ALIVE gate: the trybuild test harness,
the compile-fail and compile-pass fixture disciplines, the \texttt{.stderr}
snapshot protocol, and the certification levels from
\textsc{Alive~001} through \textsc{Crown Alive~004}.

Chapter~\ref{ch:papers} maps each of the 98~source papers to the type-law
surface it grounds: which module, which witness marker, which compile-fail
fixture encodes its core invariant.

Chapter~\ref{ch:doctrine} formalizes the manufacturing doctrine: commits as
receipts, the ALIVE gate as the certification surface, and the graduation
boundary to \texttt{wasm4pm}.

Chapter~\ref{ch:evaluation} evaluates the approach: compile-time overhead,
fixture maintenance burden, toolchain stability, and the cost of enforcing
the doctrine across 601~commits.

Chapter~\ref{ch:relatedwork} surveys related work in type-safe process
representations, liquid types, dependent types, and certified program
extraction.

Chapter~\ref{ch:conclusion} concludes with open problems and directions for
future work, including stable Rust graduation criteria and the path toward
formal mechanization in Lean~4.
